\scanpage{30}
\mspage{11}

On va supposer
\[
\boxed{\psi\sim1\qquad\text{i.e}\quad\exists\,\alpha,\beta\in\mathbb N}\quad\text{avec}\quad\boxed{\varphi^\alpha\psi=\varphi^\beta}
\]
% slanted
\margin{ce qui implique \unsure{C|\textcircled{C}|\textcircled{D}}}

Donc les relations
\[
\rho gfu\sim\rho gv\sim\rho w\sim\psi
\]
équivalent à
\[
\rho gfu\sim\rho gv\sim\rho w\sim1
\]
i.e. à
\[
\boxed{\rho gfu,\rho gv,\rho w\in M_0.}
\]

(Que $w$ ait disparu de la relation \deleted{\unsure{fin}} \unsure{ci-dessus|fondamentale} ne doit pas nous inquiéter, c'est une relation en $\deleted{\unsure{\rho,v,g,f,|f,\psi,g,q}}\underline{\rho,f,g,u,v}$, et où on peut poser $w=gv$, de sorte qu'il \unsure{en|ne} reste, comme \unsure{conditions|condition} sur ces cinq variables en question, que
\[
\underbrace{\rho gfu}_{\substack{\mathrm{def}\,\Vert\\\pi}},\ \underbrace{\rho gv}_{\substack{\Vert\,\mathrm{def}\\\pi'}}\in M_0,
\]
donc on trouve deux \unsure{élmts|éléments} $\pi,\pi'\in M_0$.

Ceci posé, on trouve les relations
\unsure{$\cdot\vert$|i)}

% archive sheet number: 29
